\chapter{Metodología}
\label{ch:metodologia}

Este capítulo presenta el marco orientado a la demanda en su totalidad. Se
organiza en ocho etapas encadenadas (\crefrange{sec:w1}{sec:w8}), resumidas en la
\cref{fig:pipeline}. El principio de diseño rector es una separación estricta entre
una capa de \emph{diagnóstico} (etapas W1--W4), que estima dónde se necesita el
transporte con independencia de dónde existe actualmente, y una capa
\emph{generativa y de validación} (W5--W8), que propone y evalúa corredores
candidatos. Todo el cómputo espacial se realiza en el sistema de referencia
proyectado canónico EPSG:6372, y todos los indicadores se agregan al \ageb, el área
geoestadística básica utilizada como unidad de análisis.

\begin{figure}[htbp]
  \centering
  \includegraphics[width=\linewidth]{pipeline.pdf}
  \caption{El marco orientado a la demanda de ocho etapas. La capa de diagnóstico
    (W1--W4) está desacoplada de la oferta existente; la capa generativa y de
    validación (W5--W8) propone corredores y los pone a prueba.}
  \label{fig:pipeline}
\end{figure}

\section{Estimación de la demanda de transporte (W1)}
\label{sec:w1}

El marco inicia estimando una superficie de demanda de transporte a partir de datos
de \emph{Nivel~1} únicamente ---variables censales, económicas y de la red vial---,
excluyendo deliberadamente cualquier insumo de transporte existente, de modo que la
demanda nunca sea función de la oferta actual.

\subsection{Generación de viajes}
Las producciones diarias $P_i$ del \ageb $i$ se estiman a partir de la población con
una amplificación por proporción joven, que refleja las mayores tasas de viaje de
las cohortes jóvenes:
\begin{equation}
  P_i = \tau \, \mathrm{pob}_i \,\bigl(1 + \gamma\, s^{\text{joven}}_i\bigr),
  \label{eq:productions}
\end{equation}
con $\tau = \num{2.5}$ viajes por persona por día y $\gamma = \num{0.10}$. Las
atracciones combinan proxies de empleo y actividad,
\begin{equation}
  A_i = w_e\, e_i + w_p\, \rho^{\text{poi}}_i a_i + w_r\, \rho^{\text{comercio}}_i a_i,
  \label{eq:attractions}
\end{equation}
con $e_i$ el proxy de empleo, $\rho^{\text{poi}}_i$ y $\rho^{\text{comercio}}_i$ las
densidades de puntos de interés y comercio, y $a_i$ el área del \ageb. El vector de
atracción se escala luego para que $\sum_i A_i = \sum_i P_i$, condición necesaria
para el modelo doblemente restringido.

\subsection{Modelo gravitacional doblemente restringido}
Los flujos origen--destino $T_{ij}$ se obtienen de un modelo de interacción espacial
doblemente restringido \parencite{wilson1971gravity},
\begin{equation}
  T_{ij} = a_i\, b_j\, P_i\, A_j\, f(d_{ij}), \qquad
  f(d) = d^{-\beta},
  \label{eq:gravity}
\end{equation}
donde $f$ es una impedancia de ley de potencia sobre la distancia euclidiana entre
centroides $d_{ij}$ y $\beta$ es el parámetro de decaimiento por distancia
(calibrado en la \cref{sec:w2}). Los factores de balanceo $a_i,b_j$ imponen las
restricciones de producción y atracción,
\begin{equation}
  a_i = \Bigl(\textstyle\sum_j b_j A_j f(d_{ij})\Bigr)^{-1}, \qquad
  b_j = \Bigl(\textstyle\sum_i a_i P_i f(d_{ij})\Bigr)^{-1},
  \label{eq:furness}
\end{equation}
resueltos por ajuste proporcional iterativo (procedimiento de Furness), alternando
ambas actualizaciones hasta que los marginales convergen (tolerancia $10^{-5}$). Los
autoflujos se anulan en la diagonal y solo se conservan los pares con
$T_{ij}\ge\num{0.5}$ en la matriz dispersa de salida.

\subsection{Superficie de demanda de transporte}
La demanda total del \ageb $D_i = \tfrac12\sum_j (T_{ij}+T_{ji})$ se pondera
finalmente por un término de propensión al transporte derivado de la propiedad
vehicular de los hogares. Con
$v_i = \mathrm{VPH\_AUTOM}_i / \max(\mathrm{VIVPAR\_HAB}_i, 1)$ la tasa vehicular,
\begin{equation}
  D^{\text{transp}}_i = D_i \,(1 - v_i),
  \label{eq:transit-demand}
\end{equation}
de modo que las zonas con alta propiedad automovilística aportan proporcionalmente
menos demanda latente de transporte. La superficie resultante $D^{\text{transp}}$ es
la cantidad dependiente que impulsa el diagnóstico de brecha en la \cref{sec:w3}.

\section{Calibración con encuesta (W2)}
\label{sec:w2}

El parámetro de decaimiento $\beta$ se calibra contra la Encuesta Origen--Destino
2022 \parencite{imeplan_eod2022}. Los \ageb{}s se unen espacialmente a las zonas de
la encuesta, los flujos modelados se agregan a pares de zonas y se elige $\beta$
para minimizar la suma de errores al cuadrado en el espacio logarítmico entre los
flujos modelados y observados,
\begin{equation}
  \beta^\star = \arg\min_{\beta}\;
  \sum_{(z,z')} \Bigl(\log \hat{T}_{zz'}(\beta) - \log T^{\text{obs}}_{zz'}\Bigr)^2 .
  \label{eq:calibration}
\end{equation}
En el universo corregido de \num{1881} \ageb{}s esto arroja
$\beta^\star = \num{1.2005}$, que supera al prior inicial $\beta = \num{2.0}$ tanto
en RMSE como en $R^2$; el valor calibrado se adopta en todo el trabajo.

\section{Capa de oferta y brecha de cobertura (W3)}
\label{sec:w3}

\subsection{Accesibilidad al transporte}
La oferta existente se mide mediante una superficie de accesibilidad de
oportunidades acumuladas construida a partir del alimentador \textsc{gtfs}. Se
construye un grafo de transporte dirigido a partir de pares consecutivos de paradas;
desde cada \ageb se identifican las paradas de abordaje dentro de \SI{400}{\metre} y
una búsqueda de ruta más corta asigna a cada destino alcanzable un tiempo de viaje
generalizado
\begin{equation}
  t = t_{\text{cam}} + t_{\text{esp}} + t_{\text{veh}}
    = \frac{d_{\text{cam}}}{s_{\text{cam}}} + \frac{h}{2} + t_{\text{veh}},
  \label{eq:traveltime}
\end{equation}
con velocidad peatonal $s_{\text{cam}} = \SI{80}{\metre\per\minute}$ e intervalo
medio $h$. La accesibilidad es el empleo total alcanzable dentro de un presupuesto
de \SI{45}{\minute},
\begin{equation}
  \mathrm{Acc}_i = \sum_{j : t_{ij}\le 45} e_j .
  \label{eq:accessibility}
\end{equation}

\subsection{Índice de brecha de cobertura}
El diagnóstico central combina demanda y oferta en un índice de brecha de cobertura,
\begin{equation}
  g_i = \frac{D^{\text{transp}}_i}{\mathrm{Acc}_i + \varepsilon},
  \qquad \varepsilon = 1,
  \label{eq:gap}
\end{equation}
normalizado por una transformación $\log(1+x)$ seguida de un escalamiento min--máx
para dar $g^{\,n}_i \in [0,1]$. Los \ageb{}s se clasifican por quintiles conjuntos:
un \ageb es de \emph{brecha alta} cuando está en los dos quintiles superiores de
demanda y en los dos inferiores de accesibilidad, de \emph{brecha baja} en el caso
espejo y de \emph{brecha media} en el resto. La \cref{fig:coverage-gap} mapea el
resultado para \zmg.

\begin{figure}[htbp]
  \centering
  \includegraphics[width=0.82\linewidth]{zmg_coverage_gap.pdf}
  \caption{Diagnóstico de brecha de cobertura para la Zona Metropolitana de
    Guadalajara. El núcleo bien servido es de brecha baja; los \ageb{}s de brecha
    alta se concentran en la periferia.}
  \label{fig:coverage-gap}
\end{figure}

\subsection{Modelo supervisado sobre la brecha de cobertura}
Para interpretar los impulsores de la demanda no atendida, un objetivo binario
$y_i = \mathbb{1}[\text{categoría-}g_i = \text{brecha alta}]$ se regresa sobre los
\num{14} indicadores Nodo--Lugar--Personas (todas las variables de oferta excluidas
para evitar circularidad) mediante Random Forest y árboles potenciados por gradiente
\parencite{niu2023rf,ke2017lightgbm}. El comportamiento del modelo se interpreta con
valores SHAP \parencite{lundberg2017shap}; población, rezago social y empleo emergen
como los impulsores dominantes del estatus de brecha alta.

\section{Priorización Nodo--Lugar--Personas (W4)}
\label{sec:w4}

Con independencia de la demanda y la oferta, se calcula una puntuación de
priorización basada en el lugar a partir de los \num{14} indicadores normalizados.
Los pesos objetivos se derivan sin intervención de expertos combinando el método
CRITIC \parencite{diakoulaki1995critic} con el método de entropía. El peso CRITIC
del indicador $j$ premia tanto la dispersión como la independencia,
\begin{equation}
  w^{\text{c}}_j \propto \sigma_j \sum_{k} \bigl(1 - r_{jk}\bigr),
  \label{eq:critic}
\end{equation}
con $\sigma_j$ la desviación estándar del indicador $j$ y $r_{jk}$ su correlación
con el indicador $k$; ambas ponderaciones se promedian en un peso de ensamble $w_j$
(\cref{fig:weights}). La puntuación de lugar, el término de equidad y la puntuación
final son
\begin{align}
  \mathrm{NPP}_i &= \textstyle\sum_j w_j\, x^{\,n}_{ij}, \label{eq:npp}\\
  \mathrm{Eq}_i  &= \tfrac12\bigl(x^{\,n,\text{marg}}_i + x^{\,n,\text{rezago}}_i\bigr),
    \label{eq:equity}\\
  F_i &= (1-\alpha)\,\mathrm{NPP}_i + \alpha\,\mathrm{Eq}_i, \qquad \alpha = 0.20.
  \label{eq:final}
\end{align}
El término de equidad es un bono aditivo explícito y separable, de modo que las
puntuaciones pueden reportarse con y sin él; un análisis de sensibilidad sobre
$\alpha \in \{0.10, 0.20, 0.30\}$ confirma que el ordenamiento es robusto
(Spearman $\ge 0.98$).

\begin{figure}[htbp]
  \centering
  \includegraphics[width=0.85\linewidth]{w4_weights.pdf}
  \caption{Pesos objetivos de ensamble CRITIC$\,\times\,$EWM para los \num{14}
    indicadores Nodo--Lugar--Personas, coloreados por dimensión.}
  \label{fig:weights}
\end{figure}

\section{Función de evaluación multiobjetivo (W5)}
\label{sec:w5}

Los corredores candidatos (de W6) y las rutas existentes (en W7) se evalúan contra
tres objetivos en competencia: ganancia ponderada por demanda $f_1$ (maximizar),
longitud de ruta $f_2$ (minimizar) y equidad media $f_3$ (maximizar):
\begin{equation}
  f_1 = \frac{\sum_{i\in S} D_i\, \kappa\, u_i}{\sum_{i\in S} D_i}, \quad
  f_2 = \ell, \quad
  f_3 = \frac{1}{|S|}\sum_{i\in S}\mathrm{Eq}_i,
  \label{eq:objectives}
\end{equation}
donde $S$ es el conjunto de \ageb{}s servidos, $u_i = g^{\,n}_i$ la fracción no
atendida, $\ell$ la longitud de ruta y $\kappa$ un factor de ganancia ($\num{0.50}$
si el candidato conecta con la red existente, $\num{0.20}$ en otro caso). Un
candidato es \emph{factible} solo si satisface todas las restricciones:
\begin{equation}
  \frac{\ell}{\ell_{\text{tramo}}} \le 1.8, \quad
  \frac{\ell\cdot 10^3}{n_{\text{paradas}}-1} \in [300,1000]\,\si{\metre}, \quad
  \textstyle\sum_{i\in S} D_i \ge 500, \quad \ell \le \SI{30}{\kilo\metre},
  \label{eq:constraints}
\end{equation}
donde $\ell_{\text{tramo}}$ es el tramo en línea recta de las anclas del candidato
(el denominador de \emph{rectitud por anclas}; véase la \cref{sec:w6}). Los
candidatos se ordenan por clasificación no dominada (Pareto) minimizando
$(-f_1, f_2, -f_3)$.

\section{Generación de corredores (W6)}
\label{sec:w6}

Los corredores nuevos se generan con el \cref{alg:w6}. Los \ageb{}s de brecha alta
se restringen a la \emph{frontera} servido/no servido ---aquellos dentro de
\SI{400}{\metre} de la red existente--- para que los corredores generados conecten
intrínsecamente con ella. Las anclas de frontera se agrupan espacialmente; dentro de
cada grupo el corredor es la trayectoria hoja-a-hoja más larga (el \emph{tronco de
diámetro}) del árbol de expansión mínima de las anclas, ensamblada a partir de
segmentos reales de la red vial. La factibilidad usa la razón de rectitud por anclas
de la \cref{eq:constraints} en lugar de una razón de rodeo por extremos, porque esta
última penaliza en exceso a un corredor de cobertura de demanda que legítimamente se
curva. A cada corredor factible se le asigna un modo por umbral de demanda (tren
ligero $\ge \num{75000}$, autobús de tránsito rápido $\ge \num{15000}$, autobús
local en otro caso, viajes por día). La \cref{fig:corridors} muestra los cuatro
corredores factibles de \zmg.

\begin{algorithm}[htbp]
\caption{Generación de corredores por anclas de frontera (W6)}
\label{alg:w6}
\begin{algorithmic}[1]
\Require superficie de brecha $g^{\,n}$, demanda $D$, red vial $G$, paradas de la red
\State $H \gets$ clase Jenks superior de $g^{\,n}$ con $D \ge 500$
  \Comment{reserva de brecha alta}
\State $H \gets \{i \in H : \operatorname{dist}(i,\text{red}) \le \SI{400}{\metre}\}$
  \Comment{frontera}
\State particionar $H$ en $k$ grupos espaciales por $k$-medias
\ForAll{grupo $C$}
  \State $T \gets$ árbol de expansión mínima de $C$ sobre $G$
  \State $\pi \gets$ trayectoria hoja-a-hoja más larga de $T$ \Comment{tronco de diámetro}
  \State construir candidato desde $\pi$; evaluar $f_1,f_2,f_3$ y restricciones (W5)
\EndFor
\State \Return candidatos factibles, ordenados por Pareto
\end{algorithmic}
\end{algorithm}

\begin{figure}[htbp]
  \centering
  \includegraphics[width=0.82\linewidth]{zmg_corridors.pdf}
  \caption{Los cuatro corredores factibles generados (G00--G03) superpuestos sobre
    la superficie de brecha de cobertura.}
  \label{fig:corridors}
\end{figure}

\section{Auditoría de rutas existentes (W7)}
\label{sec:w7}

Cada ruta \textsc{gtfs} existente se evalúa con la misma función W5 y se marca como
\emph{baja demanda} ($f_1<\num{0.2}$ y puntuación total $<\num{0.3}$),
\emph{indirecta} (razón de rodeo $>\num{1.5}$) o \emph{redundante} (traslape de
Jaccard de \ageb{}s servidos $\ge\num{0.60}$ con una ruta de mayor puntuación). Las
rutas marcadas reciben una propuesta de modificación (atajo, fusión o retiro), lo
que hace la auditoría directamente comparable con los corredores generados.

\section{Validación (W8)}
\label{sec:w8}

El marco se valida de tres formas. Una \emph{retro-prueba de exclusión} enmascara el
nivel premium de transporte del alimentador \textsc{gtfs}, recalcula la accesibilidad
y la brecha de cobertura sin él, vuelve a ejecutar el generador y mide el traslape de
trazado entre los corredores re-propuestos y las rutas enmascaradas. Una
\emph{comparación} contrasta los corredores generados con las rutas premium
existentes. Finalmente, las \emph{métricas antes/después} ---tasa de cobertura,
coeficiente de Gini de la accesibilidad y población servida por kilómetro de ruta---
cuantifican la contribución marginal de los corredores generados. El
\cref{ch:resultados} reporta las tres para \zmg y el \cref{ch:transferibilidad} las
repite para las ciudades de transferencia.
